% =============================================================================
% borrador_estudio_reunion.tex — Documento de estudio para la reunión con
% Gonzalo Gutiérrez (martes). Sincronizado con presentacion_gonzalo.pdf
% (16 láminas, versión v2 post-arbitraje simulado).
%
% NOTA: este documento REEMPLAZA la Parte B de guion_presentacion.tex, que
% describe el deck antiguo de 12 láminas con la física pre-corrección
% (equipartición sobre E0 con Nc=N). NO estudiar esa versión.
%
% Compilar (desde docs/): pdflatex borrador_estudio_reunion.tex (dos pasadas)
% =============================================================================
\documentclass[11pt]{article}

\usepackage[utf8]{inputenc}
\usepackage[T1]{fontenc}
\usepackage[spanish,es-nodecimaldot]{babel}
\usepackage[a4paper,margin=2.3cm]{geometry}
\usepackage{amsmath,amssymb}
\usepackage{booktabs}
\usepackage{enumitem}
\usepackage{xcolor}
\usepackage[colorlinks=true,linkcolor=blue!50!black]{hyperref}

\newcommand{\eps}{\varepsilon}
\newcommand{\clave}[1]{\textcolor{red!60!black}{\textbf{#1}}}
\newcommand{\decir}[1]{\begin{quote}\itshape``#1''\end{quote}}

\title{\vspace{-1.2cm}Borrador de estudio --- reunión con Gonzalo\\
\large Guion de la presentación (16 láminas, v2) + física de respaldo +
preguntas probables}
\author{Felipe Wasaff}
\date{Julio 2026}

\begin{document}
\maketitle

\section*{La historia en tres frases (memorízala)}

\begin{enumerate}[leftmargin=1.6em]
    \item La energía que un impacto ``pierde'' no se va al piso: queda
    mayoritariamente \clave{dentro de la bola}, mitad cinética
    ($\alpha_{T_B}\approx0.22$, plateau) y mitad potencial
    ($\Delta V_B\approx T_B$ por virial; el \emph{exceso} sobre $T_B$
    es la parte anarmónica).
    \item La equipartición correcta se aplica a la \clave{energía
    disipada}, no a $E_0$:
    $f_{\rm ball}=(3N-3)/(3N-3+3N_c^{\rm ef})$ con
    $N_c^{\rm ef}=25\pm8$ constante $\Rightarrow$ $f_{\rm ball}\to1$
    para $N$ grande (no $1/2$).
    \item $N_c^{\rm ef}\approx25$ es chico porque la superficie es un
    reservorio \clave{ancho pero casi elástico}: cientos de átomos se
    mueven durante el contacto ($N_c^{\rm part}\approx300$--$410$,
    medido) pero devuelven casi toda la energía al despegue.
\end{enumerate}

\noindent\textbf{Estrategia narrativa:} contar el error corregido
explícitamente. El argumento viejo ($N_c=N$, $\alpha_{\rm bola}\to1/2$)
era inconsistente; el arbitraje simulado lo detectó; la medición directa
de $N_c$ lo descartó. Presentarlo así demuestra control metodológico y se
anticipa a la objeción natural de cualquier físico estadístico ---
incluida la que haría el propio Gonzalo.

% =============================================================================
\section{Guion lámina por lámina ($\sim$45 min de exposición + discusión)}
% =============================================================================

\subsection*{Lámina 1 --- Portada (15 s)}
\decir{Te traigo el trabajo completo del rebote: la reproducción del
2012, lo nuevo que encontramos, una corrección importante que hicimos en
el camino, y cinco decisiones que necesito cerrar contigo hoy.}

\subsection*{Lámina 2 --- De dónde venimos y qué hicimos (2 min)}
\begin{itemize}[leftmargin=1.4em]
    \item Bloque derecho primero: $e<1$ dice \textsc{cuánta} energía no
    vuelve; nosotros respondemos \textsc{dónde} queda.
    \item Izquierda, sin detenerse: partición al primer apogeo, 4
    canales con cierre exacto; barrido
    $N\in\{1,2,7,13,55,79,87\}\times6$ alturas $\times$
    \clave{5 orientaciones = 210 corridas}, todo media$\,\pm\,\sigma$;
    $N_c$ \clave{medido} en el dump; estado de la bola, losa doble,
    simulación larga, adhesión.
    \item Puente: ``el experimento cabe en un video de 30 segundos.''
\end{itemize}

\subsection*{Lámina 3 --- Video 1: rebote sólido (3 min)}
\texttt{videos\_reunion/video\_N79\_v1.16\_solido.mp4} ($N{=}79$,
$h_0{=}28\,\sigma$, $v_{\rm imp}{=}1.16$; impacto en $t\approx46\,\tau$,
apogeo en $t\approx72\,\tau$; color = KE por átomo).
\begin{itemize}[leftmargin=1.4em]
    \item Deja correr el video una vez completo, sin hablar.
    \item Segunda pasada, narrada: ``mira la onda de compresión que entra
    a la bola en el impacto; los fonones que radian en la losa; y al
    separarse, la bola queda \emph{moteada}: caliente por dentro.''
    \item La figura estática (\texttt{fig\_sequence}) queda en pantalla:
    (a) caída, (b) compresión máxima, (c) separación, (d) apogeo.
    \item Remate: ``la intuición `se disipa en el suelo' falla: el piso
    se queda con menos del 15\%.''
\end{itemize}

\subsection*{Lámina 4 --- Los cuatro canales (3 min)}
\begin{itemize}[leftmargin=1.4em]
    \item ``Lo mismo, cuantificado: seis alturas, cinco orientaciones
    cada una, media y desviación.''
    \item Columna $\alpha_{T_B}$: \clave{plateau $\approx0.22$ para
    $0.9<v\lesssim1.7$} --- la fracción deja de depender de cuán fuerte
    golpeas; señal de termalización interna.
    \item Columna $\alpha_{\Delta V_B}$: canal potencial,
    $\Delta V_B\approx T_B$ por virial (sólido casi armónico); el exceso
    $\Delta V_B-T_B$, que crece con $v$, es la parte anarmónica.
    \item Cierre: $\alpha_{\Delta V_B}$ se reporta por cierre; la
    verificación con $V_B$ \clave{directo} muestra que solo falta la KE
    \clave{lateral} del CM (mediana $5\times10^{-4}E_0$;
    $<0.02\,E_0$ para $N\ge55$).
    \item Ojo: el límite superior $v\lesssim1.7$ se justifica en la
    lámina 11 (fusión); si pregunta antes, adelantar ``la bola se funde a
    $v\approx2$, lo vemos en dos láminas''.
\end{itemize}

\subsection*{Lámina 5 --- Control de calidad I: conservación (2 min)}
\begin{itemize}[leftmargin=1.4em]
    \item ``Primera pregunta obligada de cualquier simulador: ¿se
    conserva la energía? Sí: $|\Delta E|=0.023\,\eps=4.3\times10^{-4}E_0$
    en $400\,\tau$, NVE, $\Delta t=5\times10^{-3}\,\tau$.''
    \item Panel (b): $U_c$ es \clave{transitorio} (solo existe durante el
    contacto y regresa a la bola); $T_B$ y $\Delta V_B$ quedan,
    aproximadamente iguales.
    \item Anécdota técnica que da credibilidad: sin
    \texttt{pair\_modify shift yes} el cierre fallaba 2--6\% de $E_0$
    (saltos de $\pm0.016\,\eps$ por par al cruzar $r_c$). Trayectorias
    idénticas; pura contabilidad.
\end{itemize}

\subsection*{Lámina 6 --- Control de calidad II: termalización (2 min)}
\begin{itemize}[leftmargin=1.4em]
    \item Panel (a): salto de $T_{\rm ef}$ en el impacto y equilibrio
    interno en $\sim10\,\tau$ de contacto $\gg$ período vibracional ---
    \clave{esta es la condición que autoriza hablar de equipartición}.
    \item Panel (b): $T_{\rm ef}$ en el apogeo vs $v_{\rm imp}$: la bola
    queda \emph{fría} respecto de la fusión hasta $v\approx1.7$.
\end{itemize}

\subsection*{Lámina 7 --- Control de calidad III: validación contra su
Fig.~7 (3 min; tono de colaboración)}
\begin{itemize}[leftmargin=1.4em]
    \item ``Antes de extender, validamos contra ustedes. Digitalicé tu
    Fig.~7 por color de píxel, calibrada con las alturas conocidas.''
    \item El hallazgo de convención: en 2012 $h_0$ es la altura
    \clave{absoluta} del CM ($z_{\rm rest}\approx10.3\,\sigma$); nosotros
    medimos desde el tope de la losa. Misma etiqueta, caída distinta.
    \item Comparando a caída efectiva igual: \clave{0.90 vs 0.92} en el
    régimen cuasi-elástico; a alta energía diferencias 0.06--0.09,
    dentro de la dispersión entre orientaciones ($\pm0.05$--$0.14$) y
    explicables (cúmulo sin minimizar, $N$ variable 79--84, Beeman).
    \item \textsc{Nunca} decir ``error de ustedes''; decir ``lección de
    convenciones, documentada en la Sec.~3c''.
\end{itemize}

\subsection*{Lámina 8 --- El argumento que corregimos (4 min; lámina
bisagra)}
Esta lámina es la más importante de contar bien. Guion casi literal:
\decir{La primera versión del manuscrito aplicaba equipartición sobre
$E_0$ con $N_c=N$, y predecía $\alpha_{\rm bola}\to1/2$. Sometimos el
texto a un arbitraje simulado y detectó tres inconsistencias: el CM no
participa del reparto, $N_c$ tenía tres calibraciones distintas según la
sección, y la predicción $\alpha_{\rm surf}\to1/2$ no está en los datos.
La corrección: la única energía que se reparte es la \textbf{disipada},
$1-\alpha_{\rm cm}$. Eso define
$f_{\rm ball}=\alpha_{\rm bola}/(\alpha_{\rm bola}+\alpha_{\rm surf})$,
y el conteo de modos da $f_{\rm ball}=(3N-3)/(3N-3+3N_c^{\rm ef})$. Con
$N_c^{\rm ef}=25\pm8$ constante, $f_{\rm ball}\to1$: la bola grande se
queda con casi todo lo que disipa.}
\begin{itemize}[leftmargin=1.4em]
    \item Si pregunta quién fue el árbitro: arbitraje simulado interno;
    el informe está en \texttt{docs/} del repo
    (\texttt{informe\_arbitraje\_simulado}). Mostrarlo si lo pide --- es
    evidencia de método, no debilidad.
\end{itemize}

\subsection*{Lámina 9 --- El resultado central: $f_{\rm ball}(N)$ (3 min)}
\begin{itemize}[leftmargin=1.4em]
    \item Tabla: $f_{\rm ball}$ crece de $0.34$ ($N{=}2$) a $0.77$
    ($N{=}87$); ajuste global con \clave{una sola constante}
    $N_c^{\rm ef}=25\pm8$.
    \item Ser honesto sin que lo pida: ``$N_c^{\rm ef}$ es
    \emph{ajustado}, pero acotado por dos mediciones directas
    independientes --- siguiente lámina.''
    \item $N=7$ y $N=13$: régimen molecular, no termalizan; la
    \clave{excepción que confirma el mecanismo}. Acuerdo cuantitativo
    para $N\ge55$, cualitativo por debajo. (Pendiente v3: excluirlos del
    ajuste o reportar $\chi^2$/dof --- está en la lámina 15.)
\end{itemize}

\subsection*{Lámina 10 --- Las tres escalas de $N_c$ (3 min)}
\begin{itemize}[leftmargin=1.4em]
    \item La jerarquía medida en el dump ($v\approx1.47$):
    \[ N_c^{\rm ef}\approx25 \;\ll\; N_c^{\rm geom}\approx45\text{--}90
       \;\ll\; N_c^{\rm part}\approx300\text{--}410. \]
    \item La frase que resume el paper: \decir{la superficie es un
    reservorio ancho pero casi elástico: cientos de átomos se mueven
    transitoriamente, pero devuelven la energía al despegue; el sumidero
    neto equivale a solo $\sim$25 osciladores.}
    \item El viejo $N_c\approx290$ coincidía \clave{por accidente} con
    $N_c^{\rm part}$: número correcto, balance equivocado.
\end{itemize}

\subsection*{Lámina 11 --- Video 2 + fusión (3 min)}
\texttt{videos\_reunion/video\_N79\_v2.0\_fusion.mp4} ($v_{\rm imp}
\approx2.0$; impacto en $t\approx78\,\tau$, apogeo $t\approx124\,\tau$).
\begin{itemize}[leftmargin=1.4em]
    \item ``A $v\approx2$ el rebote es cualitativamente distinto: mira
    cómo la bola pierde la forma.''
    \item $g(r)$: a $v=1.47$ picos FCC nítidos (contraste $\approx11$);
    a $v=2.03$ colapsada (contraste $\approx2$) --- \clave{fundida},
    consistente con la fusión que el 2012 reportó.
    \item Matiz importante: a $v=2.03$, $T_{\rm ef}=0.28$ queda
    \emph{bajo} el rango de fusión ($0.30$--$0.35$); el criterio decisivo
    es $g(r)$, no $T_{\rm ef}$.
    \item Consecuencia: plateau y conteo de modos se restringen a
    $v\lesssim1.7$.
\end{itemize}

\subsection*{Lámina 12 --- Tamaño finito de la losa (2 min)}
\begin{itemize}[leftmargin=1.4em]
    \item ``El árbitro tenía razón'': con losa $16\times16\times12$,
    $\alpha_{\rm surf}$ sube $0.083\to0.143$ ($\sim0.06\,E_0$ de sesgo
    por fonones reflejados) y $f_{\rm ball}$ baja $0.87\to0.75$.
    \item \clave{La conclusión no cambia}: $f_{\rm ball}=0.75$ aun con
    losa doble. Sesgo declarado junto a los datos.
    \item Punto débil conocido: es \emph{un} caso de control; un segundo
    (otro $N$ u otra altura) son $\sim$30 min de cómputo si lo pide.
\end{itemize}

\subsection*{Lámina 13 --- $h^*=T_B/(Mg)$ (2 min)}
\begin{itemize}[leftmargin=1.4em]
    \item Primer rebote: $h_1/h^*=0.96$ --- la escala funciona
    (predicción de su 2012).
    \item No es atractor: la razón cae monótonamente; atrapamiento tras
    23 rebotes ($t\approx660\,\tau$).
    \item El ``colapso'' final a través de la red: artefacto del
    potencial repulsivo puro (sin barrera lateral), documentado como
    ejercicio de análisis crítico.
\end{itemize}

\subsection*{Lámina 14 --- Adhesión (2 min)}
\begin{itemize}[leftmargin=1.4em]
    \item Con $c_{ps}=1$: 24/24 casos adheridos hasta $v=2$.
    \item $E_{\rm adh}\propto N^{0.73}$ ($\approx$ área de contacto
    $N^{2/3}$); masa $\propto N$ $\Rightarrow$ los grandes escapan más
    fácil ($h_0^*\propto N^{-0.27}$) --- conexión con \emph{soft
    landing}.
    \item Umbral extrapolado $h_0^*\sim10^2$--$10^3\,\sigma$: solo orden
    de magnitud (a esas $v$ la bola se funde).
\end{itemize}

\subsection*{Lámina 15 --- Estado editorial (2 min)}
\begin{itemize}[leftmargin=1.4em]
    \item Hecho: 5 puntos mayores del arbitraje resueltos \clave{con
    simulaciones nuevas}; manuscrito v2 compila (9 págs), protocolo RBEF
    cumplido (bilingüe, \texttt{bibitem} en orden, sin TikZ).
    \item Falta ($\approx$1 h, nada es cómputo): su ORCID (bloqueante),
    frase sobre $N=7$ en el ajuste, criterio de plateau explícito,
    recorte del resumen, ``surprising'' del abstract EN, párrafo de
    público objetivo.
\end{itemize}

\subsection*{Lámina 16 --- Decisiones (5 min; cierre operativo)}
Recorrer las 5 y \clave{anotar las respuestas}:
\begin{enumerate}[leftmargin=1.6em]
    \item \textbf{Atribución}: plateau $\alpha_{T_B}$, fórmula de
    equipartición sobre disipada, $h^*$ --- ¿contribución nueva o notas
    previas suyas? Si algo es suyo, se cita; no es conflicto, es
    limpieza de autoría.
    \item \textbf{ORCID} de Gonzalo (único bloqueante logístico).
    \item \textbf{Idioma} del envío (español aceptado por RBEF).
    \item \textbf{Orden de autores} / correspondencia (hoy: Wasaff,
    Gutiérrez).
    \item \textbf{Flujo}: ¿arXiv \texttt{physics.ed-ph} simultáneo?
    ¿Versión extendida (doble columna) para una revista de física
    computacional después?
\end{enumerate}

% =============================================================================
\section{Física de respaldo (versión v2 --- estudiar esta, no la vieja)}
% =============================================================================

\subsection{Los cuatro canales y el cierre}
Al primer apogeo, $E_0=KE$ de impacto se reparte en
\[ \alpha_{\rm cm}=e_1^2=\frac{h_1}{h_0},\qquad
   \alpha_{T_B}=\frac{T_B}{E_0},\qquad
   \alpha_{\Delta V_B}=\frac{\Delta V_B}{E_0},\qquad
   \alpha_{\rm surf}=\frac{\Delta E_{\rm surf}}{E_0}, \]
con $\alpha_{\rm bola}\equiv\alpha_{T_B}+\alpha_{\Delta V_B}$. En las
tablas $\alpha_{\Delta V_B}$ se obtiene por cierre ($\sum\alpha_i=1$);
la verificación independiente con $V_B$ medido muestra que el único
término ausente es la KE \emph{lateral} del CM en el apogeo: mediana
$5\times10^{-4}E_0$, $<0.02\,E_0$ para $N\ge55$ (hasta $0.06\,E_0$ en
$N=2,13$ a baja $v$ --- refuerza el carácter molecular de esos casos).
En el caso canónico ($N=79$, $h_0=28.15$) el cierre directo es
$<10^{-4}$.

\subsection{El plateau}
Para $0.9<v_{\rm imp}\lesssim1.7\,\sigma/\tau$ (y $N\ge55$),
$\alpha_{T_B}$ es constante $\approx0.22$: la \emph{fracción} de energía
que termina como KE interna no depende de la violencia del impacto.
Interpretación: el contacto ($\sim10\,\tau$) dura decenas de períodos
vibracionales y la bola termaliza internamente; el reparto queda fijado
por el conteo de modos, no por $E_0$. Por debajo de $v\approx0.9$:
régimen cuasi-elástico. Por encima de $v\approx1.7$: fusión.

\subsection{Equipartición sobre la energía disipada (la corrección)}
El CM no participa del reparto: es un modo macroscópico que \emph{se
lleva} $\alpha_{\rm cm}$. Lo que se reparte entre modos es lo disipado,
$1-\alpha_{\rm cm}$, entre los $3N-3$ modos internos de la bola y los
$3N_c^{\rm ef}$ grados efectivos del sumidero neto de la superficie:
\[ f_{\rm ball}\equiv
   \frac{\alpha_{\rm bola}}{\alpha_{\rm bola}+\alpha_{\rm surf}}
   =\frac{3N-3}{3N-3+3N_c^{\rm ef}},\qquad N_c^{\rm ef}=25\pm8. \]
Límites: $f_{\rm ball}\to1$ para $N\to\infty$ (la superficie neta no
crece con $N$); validez declarada: bola fría, sin adhesión, superficie
cristalina casi elástica, un solo impacto. No se extrapola a pelotas
macroscópicas.

\subsection{Las tres escalas de $N_c$ (todas con significado físico)}
\begin{center}
\begin{tabular}{llll}
\toprule
Escala & Valor & Cómo se obtiene & Qué significa \\
\midrule
$N_c^{\rm ef}$   & $25\pm8$    & ajuste global de $f_{\rm ball}(N)$ & sumidero \emph{neto} \\
$N_c^{\rm geom}$ & 45--90      & huella geométrica en el dump       & contacto directo \\
$N_c^{\rm part}$ & 300--410    & átomos con KE $>$ umbral           & participación transitoria \\
\bottomrule
\end{tabular}
\end{center}
La jerarquía $N_c^{\rm ef}\ll N_c^{\rm geom}\ll N_c^{\rm part}$ es el
contenido físico: reservorio ancho, casi elástico.

\subsection{$T_B\approx\Delta V_B$: el virial}
En un sólido casi armónico, $\langle KE\rangle=\langle\Delta
PE\rangle$ por modo. Que los dos canales internos salgan iguales
(visible en \texttt{fig\_conservacion}b y en las tablas) es evidencia
independiente de termalización interna --- no un supuesto.

\subsection{Números de contexto}
Unidades de argón: $\eps\approx0.0104$~eV, $\sigma\approx3.4$~\AA,
$\tau\approx2.15$~ps. Contacto efectivo $z_0\approx11.6$--$11.8\,\sigma$;
$g=0.026\,\sigma/\tau^2$; $\Delta t=0.005\,\tau$. Shells FCC completos:
$N=13,55,79,87$ ($N{=}79$: $r<2.45\,\sigma$, cluster canónico). El
``$N=83$'' antiguo era artefacto de \texttt{create\_atoms} con centro
fuera de sitio de red.

% =============================================================================
\section{Preguntas probables y respuestas (resumen operativo)}
% =============================================================================
Versión condensada; el detalle completo está en
\texttt{docs/plan\_reunion\_gonzalo.md} (imprimirlo también).

\begin{description}[leftmargin=1.2em,style=nextline]
    \item[¿Por qué sus $e_1$ no coinciden con mi Fig.~7?] Convención de
    $h_0$ (absoluta vs desde el tope de la losa). A caída efectiva
    igual: 0.90 vs 0.92; resto dentro de la dispersión entre
    orientaciones.
    \item[El 2012 ya mostraba que $U_c$ domina. ¿Qué agregan?] El 2012
    describe el balance \emph{durante} el contacto; nosotros la
    partición \emph{al apogeo}, cuando los canales ya se separaron.
    $U_c$ es transitorio y regresa a la bola. Nuevo: 4 $\alpha$ con
    cierre, plateau, equipartición sobre disipada, $N_c$ medido, $h^*$,
    adhesión.
    \item[¿$N_c^{\rm ef}$ es medido o ajustado?] Ajustado --- un
    parámetro global sobre $f_{\rm ball}(N)$, acotado por dos mediciones
    directas ($N_c^{\rm geom}$, $N_c^{\rm part}$). No sobrevender.
    \item[¿Equipartición sin equilibrio?] Aproximada, en el subsistema
    que sí termaliza durante el contacto. Dos evidencias:
    $T_B\approx\Delta V_B$ y el éxito del conteo de modos para
    $N\ge55$. Los fallos ($N=7,13$) son los que no termalizan.
    \item[¿$f_{\rm ball}\to1$ vale en el límite macroscópico?] No se
    extrapola: régimen declarado (fría, sin adhesión, cristalina, un
    impacto). El punto es que el reservorio interno crece como $3N-3$ y
    el neto de la superficie no crece.
    \item[¿La losa no es muy chica?] Correcto y cuantificado: sesgo
    $\sim0.06\,E_0$ en $\alpha_{\rm surf}$; $f_{\rm ball}$ $0.87\to0.75$
    con losa doble; conclusión intacta.
    \item[¿La bola sigue sólida?] Sí hasta $v\approx1.7$ (criterio
    $g(r)$, contraste $\approx11$ vs $\approx2$ fundida). A $v=2.03$,
    $T_{\rm ef}=0.28<T_{\rm fus}$: el criterio decisivo es $g(r)$.
    \item[¿Barras de error?] 210 corridas (5 orientaciones); todo
    media$\,\pm\,\sigma$. La dispersión entre orientaciones es grande y
    es parte del mensaje (sin promediar, el ruido parece física).
    \item[¿Rotación / momento angular?] No medido; declarado futuro
    (\texttt{compute angmom} es una línea). Está dentro de $T_B$, no
    afecta el cierre. La traslación \emph{lateral} del CM sí queda fuera
    de los 4 canales: medida, mediana $5\times10^{-4}E_0$ y
    $<0.02\,E_0$ para $N\ge55$; declarada en el manuscrito.
    \item[¿Temperatura inicial finita?] Futuro. $T=0$ aísla el
    mecanismo; el plateau debería sobrevivir mientras
    $T\ll T_{\rm fus}$.
    \item[¿Por qué RBEF?] Ángulo pedagógico genuino (corre en un
    notebook, actividades, repo público), sin costos, acceso abierto.
    La versión de investigación existe para una segunda revista ---
    decisión 5.
\end{description}

% =============================================================================
\section*{Chuleta de números}
% =============================================================================
\begin{center}
\small
\begin{tabular}{ll}
\toprule
Plateau & $\alpha_{T_B}\approx0.22$ ($0.9<v\lesssim1.7$, $N\ge55$) \\
Equipartición & $f_{\rm ball}=(3N-3)/(3N-3+3N_c^{\rm ef})$, $N_c^{\rm ef}=25\pm8$ \\
Tres escalas & $25\ll45$--$90\ll300$--$410$ (ef / geom / part) \\
$f_{\rm ball}$ medidos & 0.34 (N2), 0.55 (N7), 0.38 (N13), 0.49 (N55), 0.70 (N79), 0.77 (N87) \\
Conservación & $|\Delta E|=4.3\times10^{-4}E_0$ en $400\,\tau$; KE lateral CM: med.\ $5\times10^{-4}E_0$, $<0.02\,E_0$ ($N\ge55$) \\
Validación 2012 & $e_1$: 0.90 (2012) vs 0.92 (2026) cuasi-elástico; $z_{\rm rest}\approx10.3\,\sigma$ \\
Losa doble & $\alpha_{\rm surf}$: $0.083\to0.143$; $f_{\rm ball}$: $0.87\to0.75$ \\
Fusión & sólida $v\le1.47$ (contraste 11); fundida $v=2.03$ (contraste 2); $T_{\rm fus}^*\approx0.30$--$0.35$ \\
$h^*$ & $h_1/h^*=0.96$ primer rebote; atrapamiento tras 23 rebotes ($t\approx660\,\tau$) \\
Adhesión & 24/24 pegados; $E_{\rm adh}\propto N^{0.73}$; $h_0^*\propto N^{-0.27}$ \\
Estadística & 7 tamaños $\times$ 6 alturas $\times$ 5 orientaciones = 210 corridas \\
Videos & sólido: impacto $46\,\tau$, apogeo $72\,\tau$ / fusión: impacto $78\,\tau$, apogeo $124\,\tau$ \\
\bottomrule
\end{tabular}
\end{center}

\section*{Checklist del día anterior}
\begin{itemize}[leftmargin=1.4em]
    \item[$\square$] Probar \texttt{presentacion\_gonzalo.pdf} y los 2
    videos en el notebook/proyector.
    \item[$\square$] Imprimir \texttt{borrador\_rbef\_v2.pdf} (9 págs) +
    este borrador + \texttt{plan\_reunion\_gonzalo.md}.
    \item[$\square$] \texttt{git push} --- repo al día.
    \item[$\square$] A mano: \texttt{tabla\_orient.txt},
    \texttt{fig\_conservacion.pdf}, \texttt{fig\_termalizacion.pdf},
    informe de arbitraje simulado.
\end{itemize}

\end{document}
